\renewcommand{\bibname}{Quellen und Bandzuordnung}
\begin{thebibliography}{9}
\addcontentsline{toc}{chapter}{Quellen und Bandzuordnung}
\footnotesize
\setlength{\itemsep}{4pt}
\setlength{\parskip}{0pt}
\interlinepenalty=10000

\bibitem{project}
Martin Kunze: \emph{Die Grundlagen der Mathematik}.
Band~5: funktionale Eindeutigkeit und Funktionen.
Band~4: totale Relationen.
Band~3: Mengenoperationen, geordnete Paare, kartesische Produkte,
Aussonderung und Potenzmengen.
Format nach den Leitfäden zu Band~3, Band~4, Band~28 und Band~37.
Der vorliegende Leitfaden bleibt bei allgemeinen Funktionen
auf Mengen und ihren unmittelbaren mengentheoretischen Folgerungen.

\bibitem{meseguer}
José Meseguer: \emph{Set Theory and Algebra in Computer Science:
A Gentle Introduction to Mathematical Modeling}.
University of Illinois at Urbana-Champaign, Fassung vom
29.~August 2024.
Abschnitt~5.1, Druckseiten~29--30: Funktionen und Funktionenmengen.
Abschnitt~5.3, Druckseiten~32--33: Bilder, Urbilder,
Zielmengenwechsel und Einschränkungen.
Abschnitt~5.4, Lemma~2 auf Druckseiten~33--34:
Assoziativität der Komposition;
Aufgabe~29 auf Druckseite~36: paarwertige Funktionen.
Abschnitt~5.5, Druckseiten~37--38: Funktionen auf Funktionenmengen,
Auswertung und Currying.
\url{https://courses.physics.illinois.edu/cs476/fa2024/readings/STACS.pdf}.

\bibitem{welch}
P.~D. Welch: \emph{Set Theory}.
University of Bristol, Fassung vom 15.~September 2026.
Abschnitt~1.3.2, Druckseiten~11--12:
Definitionen~1.24--1.27 zu Definitionsbereich, Funktionsgraph,
Einschränkung und Funktionenmenge.
Aufgabe~1.21(ii), Druckseite~12:
Vereinigung von Funktionen, die auf gemeinsamen
Definitionsbereichen übereinstimmen.
\url{https://people.maths.bris.ac.uk/~mapdw/ST/current-settheory.pdf}.

\bibitem{axler}
Sheldon Axler: \emph{Measure, Integration \& Real Analysis}.
Springer, 2020; aktualisiertes Autoren-PDF vom 12.~Juni 2026.
Abschnitt~2B, Unterabschnitt \emph{Inverse Images},
Druckseiten~29--30: Definition~2.31 sowie Sätze~2.33 und~2.34
zu Urbildern von Mengenoperationen und zur Komposition.
Die hier herangezogenen Aussagen setzen ausschließlich
Mengen und allgemeine Funktionen voraus.
\url{https://measure.axler.net/MIRA.pdf}.

\end{thebibliography}

\noindent\textbf{Einordnung der Zusammenstellung.}
Dieser Leitfaden ist eine ausgewählte Referenzliste von Definitionen,
Charakterisierungen und elementaren Folgerungen.
Er beansprucht weder Vollständigkeit aller bekannten Aussagen
noch historische Erstentdeckungen.
Die Quellen dienen der Begriffs-, Notations- und Ergebniskontrolle;
nicht jede hier formulierte Zusammenfassung steht dort
in genau dieser Gestalt.
Die Aussagen werden in der Graphkonvention des Hauptbands
und einschließlich ihrer leeren Randfälle formuliert.
Insbesondere werden bei indizierten Durchschnitten die jeweiligen
Grundmengen beachtet und keine Auswahlannahmen ergänzt.
